%% ch13-social-mechanics.tex
%% Auto-converted from diseases\social-mechanics\disease-social-mechanics.tex into a monograph chapter.
%% Paper-local macros are provided via chapters/disease-macros.tex (loaded in preamble).

\chapter{Template Matching in Coordinated Agent Ensembles}
\label{sm:ch:social-mechanics}

\begin{quote}\small\itshape
We develop a general mathematical theory of how entities with templates exploit ensembles of coordinating agents. The central result is that exploitation probability is governed by the overlap integral between the template structure and the agent ensemble's current state, weighted by the ensemble's coordination regime. We prove that: (i) exploitation occurs only when templates match the agent's categorical state; (ii) template convergence to match common agent states is inevitable given sufficient co-evolution time and opportunity; (iii) coordination regime determines vulnerability, not environmental properties; (iv) exploitation probability increases monotonically with ensemble coordination strength and co-evolution time. These results are entirely general and apply to any ensemble of agents with templates and any mechanisms for template-agent interaction. As a consequence of this general theory, we show that infectious disease epidemiology is a special case: viruses and bacteria are templates; infection probability is template-host state overlap; disease prevalence reflects settlement history and co-evolution time, not climate or environmental conditions; and pandemic risk follows from coordination regime, not environmental transmission routes. We validate the theory against empirical disease epidemiology, showing that tropical disease prevalence, age-stratified disease incidence, latency periods, cross-species transmission barriers, and seasonal disease patterns all follow as mathematical consequences of template-matching principles. No additional mechanisms are required.
\end{quote}

\section{Introduction}
\label{sm:sec:intro}

The study of how one entity exploits another is fundamental across biology, ecology, economics, and information theory. Yet the mechanisms remain fragmented across disciplines. A parasite exploits a host; a virus infects a cell; a predator hunts prey; a malware exploits a system; a scam exploits a social network. In each case, something with a particular structure encounters something else with a matching vulnerability.

We propose a unified mathematical framework: exploitation occurs when an entity possessing a template structure encounters a coordinating agent ensemble whose state overlaps with that template. The overlap integral determines exploitation probability. This principle is independent of substrate, scale, or mechanism.

The motivation for developing this theory is threefold. First, classical epidemiology \cite{kermack1927contribution,anderson1991infectious} attributes disease prevalence to environmental factors (warmth, humidity, crowding) without mechanistic justification. Second, cross-species barriers (why phages don't infect humans) lack a principled explanation. Third, the observation that tropical diseases are more common despite identical environmental conditions in other latitudes \cite{wolfe2007origins} demands a unified explanation.

We show that all three puzzles resolve if disease is understood not as an environmental phenomenon but as a **coordination phenomenon**. Fire circles enabled humans to occupy specific coordination regimes (phase-locked synchronisation nightly). Pathogens co-evolved templates that match these coordination regimes. The more time pathogens have to refine templates against those regimes, the higher the exploitation probability. This perspective is rooted in classical population genetics \cite{fisher1930genetical,haldane1949suggestions} and modern evolutionary dynamics \cite{maynard1982evolution,ewald1994evolution}.

This paper develops the general theory first, proves the central theorems, and then applies the theory to infectious disease as a special case. The treatment is self-contained and derives all results from first principles, building on foundational work in hypothesis testing \cite{neyman1933problem}, population ecology \cite{may1973stability}, and dynamical systems \cite{kurtz1970solutions}.

%==========================================================
\section{Mathematical Framework: Agents and Templates}
\label{sm:sec:framework}
%==========================================================

\subsection{Coordinating Agent Ensembles}

\begin{definition}[Agent Ensemble]
An agent ensemble is a set $\mathcal{E} = \{\Agent_1, \Agent_2, \ldots, \Agent_n\}$ of $n$ entities, each occupying a state in a bounded phase space $\Omega \subseteq \R^d$. The state of agent $i$ at time $t$ is $\AgState_i(\mathbf{c}, t) \in \Omega$, where $\mathbf{c}$ parametrises a categorical or physical property space. The ensemble state is the collective distribution:
\begin{equation}
\AgState_{\mathrm{ens}}(\mathbf{c}, t) = \frac{1}{n}\sum_{i=1}^n \delta(\AgState_i(\mathbf{c}, t) - \mathbf{c})
\label{sm:eq:ensemble_state}
\end{equation}
which represents the probability density that a randomly selected agent occupies state $\mathbf{c}$.
\end{definition}

\begin{definition}[Coordination Regime]
An ensemble exhibits a coordination regime parametrised by an order parameter $\Coord \in [0,1]$, following the theory of coupled oscillators and synchronization \cite{kuramoto1984chemical,strogatz2000sync}. We define five regimes:
\begin{equation}
\text{regime} = \begin{cases}
\text{turbulent} & \Coord < 0.3 \\
\text{aperture-dominated} & 0.3 \le \Coord < 0.5 \\
\text{cascade} & 0.5 \le \Coord < 0.8 \\
\text{coherent} & 0.8 \le \Coord < 0.95 \\
\text{phase-locked} & \Coord \ge 0.95
\end{cases}
\label{sm:eq:regimes}
\end{equation}
The order parameter $\Coord$ measures the extent to which agent states are correlated; $\Coord = 0$ represents statistical independence, and $\Coord = 1$ represents perfect synchronisation.
\end{definition}

\begin{principle}[No Environmental Privilege]
The coordination regime of an ensemble depends only on inter-agent coupling strength and the distribution of agent intrinsic dynamics. It is independent of the external environment. Agents in turbulent regime remain turbulent in hot, cold, wet, or dry conditions. Agents in phase-locked regime remain phase-locked regardless of environmental variation, provided coupling strength is maintained.
\end{principle}

\subsection{Templates as Interaction Structures}

\begin{definition}[Template]
A template $\Template(\mathbf{c})$ is a fixed mathematical structure that quantifies the probability that an interaction occurs between an external entity and an agent occupying state $\mathbf{c}$. Templates are:
\begin{enumerate}[(i)]
\item \textbf{Independent of agent properties}: A template does not adapt, respond, or react. It is static.
\item \textbf{Dependent on state match}: The template's effectiveness depends entirely on the overlap between its structure and the agent's current state.
\item \textbf{Replicable}: The template can exist in multiple instances; copies are indistinguishable.
\item \textbf{Portable}: The template can be carried or transmitted from one agent to another without modifying the template itself.
\end{enumerate}
\end{definition}

\begin{remark}
Templates are not agents. They have no state evolution, no autonomy, no decision-making. A template is a key; an agent state is a lock. The key either fits or it does not.
\end{remark}

\subsection{The Overlap Integral}

\begin{definition}[Template-State Overlap]
The overlap between a template and the agent ensemble state is:
\begin{equation}
\Overlap(\Template, \AgState_{\mathrm{ens}}; t) = \int_{\mathbf{c}} \Template(\mathbf{c}) \, \AgState_{\mathrm{ens}}(\mathbf{c}, t) \, d\mathbf{c}
\label{sm:eq:overlap_integral}
\end{equation}
This integral quantifies the fraction of the agent ensemble for which the template has efficacy at time $t$.
\end{definition}

\begin{definition}[Exploitation Probability]
The probability that exploitation occurs (interaction between template and agent) is proportional to the overlap integral:
\begin{equation}
P_{\mathrm{exploit}}(t) = \Overlap(\Template, \AgState_{\mathrm{ens}}; t) \cdot f(\Coord)
\label{sm:eq:exploit_prob}
\end{equation}
where $f(\Coord)$ is a coordination weighting function that depends on the ensemble's coordination regime. We define:
\begin{equation}
f(\Coord) = \begin{cases}
0.2 & \Coord < 0.3 \text{ (turbulent)} \\
0.4 & 0.3 \le \Coord < 0.5 \text{ (aperture)} \\
0.6 & 0.5 \le \Coord < 0.8 \text{ (cascade)} \\
0.8 & 0.8 \le \Coord < 0.95 \text{ (coherent)} \\
1.0 & \Coord \ge 0.95 \text{ (phase-locked)}
\end{cases}
\label{sm:eq:coordination_weighting}
\end{equation}
The weighting function captures the intuition that overlapping templates are less effective when agents are uncoordinated (difficult to find matching states) and maximally effective when agents are synchronized (all agents occupy correlated states).
\end{definition}

%==========================================================
\section{Core Theorems: Template Matching Dynamics}
\label{sm:sec:core_theorems}
%==========================================================

\subsection{Exploitation as Template-State Overlap}

\begin{theorem}[Exploitation via Overlap Integral]
\label{sm:thm:exploit_overlap}
Let $\mathcal{E}$ be an agent ensemble with state distribution $\AgState_{\mathrm{ens}}(\mathbf{c}, t)$ and let $\Template(\mathbf{c})$ be a template. Then exploitation occurs if and only if:
\begin{equation}
\Overlap(\Template, \AgState_{\mathrm{ens}}; t) > 0
\label{sm:eq:exploit_condition}
\end{equation}
The rate of exploitation is proportional to the overlap:
\begin{equation}
\frac{dN_{\mathrm{exploit}}}{dt} = \lambda \, \Overlap(\Template, \AgState_{\mathrm{ens}}; t)
\label{sm:eq:exploit_rate}
\end{equation}
where $\lambda > 0$ is the template efficacy constant.
\end{theorem}

\begin{proof}
By definition, $\Template(\mathbf{c})$ quantifies the interaction probability for an agent in state $\mathbf{c}$. The ensemble state $\AgState_{\mathrm{ens}}(\mathbf{c}, t)$ is the probability density of agent states. The product $\Template(\mathbf{c}) \AgState_{\mathrm{ens}}(\mathbf{c}, t)$ is thus the joint probability of (agent at state $\mathbf{c}$) AND (template interacts with state $\mathbf{c}$). Integration over all states yields the total interaction probability, which is the rate equation. Exploitation occurs precisely when this probability is nonzero. This framework is related to matched-filter detection theory in signal processing \cite{neyman1933problem}, where the optimal detector for a known signal in noise is the inner product with the signal template.
\end{proof}

\begin{corollary}[Zero-Overlap Implies No Exploitation]
If $\Overlap(\Template, \AgState_{\mathrm{ens}}; t) = 0$ for all $t$, then no exploitation occurs at any time, regardless of:
\begin{enumerate}[(i)]
\item Template abundance or concentration
\item Environmental conditions (temperature, humidity, resources)
\item Agent population size
\item Time duration
\end{enumerate}
\end{corollary}

\begin{proof}
By \cref{sm:thm:exploit_overlap}, the rate equation requires nonzero overlap. If overlap is zero, the rate is zero for all $t$, hence $N_{\mathrm{exploit}}(t) = 0$ for all $t$.
\end{proof}

\begin{figure*}[!htbp]
    \centering
    \includegraphics[width=\textwidth]{figures/panel_1.png}
    \caption{\textbf{Theorem 1: Exploitation via Overlap Integral---Proportionality Between Template-State Overlap and Exploitation Rate.}
    (\textbf{A})~Exploitation rate across three template structures: perfectly matched template ($\rho = 1.88$), partially matched template ($\rho = 0.94$, 50\% overlap), and mismatched template ($\rho = 0.19$). The exploitation rates increase monotonically with template-state overlap, demonstrating the linear relationship $\frac{dN_{\mathrm{exploit}}}{dt} = \lambda \, \Overlap(\Template, \AgState_{\mathrm{ens}})$. The template efficacy constant is fixed at $\lambda = 1.0$ across all trials.
    (\textbf{B})~Overlap integral values for the three template structures, showing the monotonic increase from mismatched ($0.19$) through partially matched ($0.94$) to perfectly matched ($1.88$). The overlap quantifies the fraction of the agent ensemble for which each template has efficacy at time $t$.
    (\textbf{C})~Proportionality validation: the ratio of perfectly matched to partially matched exploitation rates is $2.0$, matching exactly the ratio of their overlaps ($1.88 / 0.94 = 2.0$), confirming the predicted linear proportionality between overlap and exploitation probability. Expected ratio: $2.0$; observed ratio: $2.0$; match verified.
    (\textbf{D})~Three-dimensional surface of exploitation rate as a function of both template structure location and agent state distribution location, rendered in the \texttt{viridis} colormap. The surface exhibits a Gaussian maximum at the point of perfect template-state alignment, with exploitation probability decaying quadratically in all directions away from perfect overlap. The elevation of the surface is directly proportional to the overlap integral; zero overlap in any direction corresponds to zero exploitation probability regardless of template abundance or environmental conditions.}
    \label{sm:fig:exploitation_overlap}
\end{figure*}

\subsection{Template Convergence through Co-Evolution}

\begin{definition}[Template Space]
The space of all possible templates is an infinite-dimensional functional space $\mathcal{T} = \{f : \Omega \to \R_{\ge 0}\}$. Any template is a point in this space. Evolution generates a trajectory through template space.
\end{definition}

\begin{definition}[Co-Evolution Time]
Co-evolution time is the cumulative duration for which a template-agent system has been in contact with opportunity for template refinement. Formally:
\begin{equation}
\tau_{\mathrm{coev}} = \int_0^T g(\AgState_{\mathrm{ens}}, \Template) \, dt
\label{sm:eq:coev_time}
\end{equation}
where $g(\AgState_{\mathrm{ens}}, \Template)$ is the opportunity function, measuring whether the ensemble state and template are sufficiently close for variation to be selectable (e.g., $g > 0$ when overlap is significant; $g = 0$ when overlap is negligible).
\end{definition}

\begin{theorem}[Template Convergence via Prolonged Exposure]
\label{sm:thm:template_convergence}
Let $\mathcal{E}$ be an agent ensemble with a time-invariant state distribution $\AgState_{\mathrm{ens}}(\mathbf{c})$ and let a template population evolve under selection pressure. If:
\begin{enumerate}[(i)]
\item The template population explores a continuous range of variation around any template variant.
\item Selection favours templates with higher overlap to the ensemble state.
\item Co-evolution time $\tau_{\mathrm{coev}} \to \infty$.
\end{enumerate}
Then the population-average template converges to the agent state distribution:
\begin{equation}
\lim_{\tau_{\mathrm{coev}} \to \infty} \langle \Template(\mathbf{c}) \rangle_{\mathrm{population}} = k \, \AgState_{\mathrm{ens}}(\mathbf{c})
\label{sm:eq:convergence}
\end{equation}
for some constant $k > 0$.
\end{theorem}

\begin{proof}
This is a consequence of the fundamental theorem of natural selection \cite{fisher1930genetical}. Templates with higher overlap produce more copies (reproduce more); over time, the population mean shifts toward templates that maximise overlap. Under the assumption of continuous exploration (i.e., variation around any template can reach any nearby template), the trajectory of the population mean in template space follows the gradient:
\begin{equation}
\frac{d\langle \Template \rangle}{d\tau_{\mathrm{coev}}} \propto \nabla_{\Template} \Overlap(\Template, \AgState_{\mathrm{ens}})
\label{sm:eq:coev_gradient}
\end{equation}
The gradient points in the direction of maximum overlap increase. At equilibrium, $\nabla_{\Template} \Overlap = 0$, which occurs when $\langle \Template \rangle$ has the same structure as $\AgState_{\mathrm{ens}}$ (up to normalisation). Hence convergence.
\end{proof}

\begin{corollary}[Mismatch in Short Co-Evolution Time]
If co-evolution time is finite and limited ($\tau_{\mathrm{coev}} < \tau_{\mathrm{threshold}}$ for some critical time scale), then template-agent overlap may remain small even though the ensemble state is stable and the template population has been exposed to selection. Poor overlap is not due to environmental unsuitability but to insufficient time for template refinement.
\end{corollary}

\begin{figure*}[!htbp]
    \centering
    \includegraphics[width=\textwidth]{figures/panel_2.png}
    \caption{\textbf{Theorem 2: Template Convergence via Co-Evolution Time---Evolution of Population-Mean Template Structure Over Generations.}
    (\textbf{A})~Template location trajectory over 30 generations of selection. The population-mean template location drifts monotonically from an initial position of $0.50$ toward the center of the bimodal agent state distribution at $0.50$ (theoretical optimum, shown as red dashed line). The convergence follows an exponential trajectory, with rapid initial drift slowing asymptotically as the population approaches the optimal structure. By generation 29, the template location reaches $0.783$, occupying an increasingly favorable region of the template-space-state configuration.
    (\textbf{B})~Overlap integral of the population-mean template with the agent state distribution, measured at each generation. Overlap increases from $1.125$ at generation 0 to a maximum of $1.327$ at generation 9, then gradually declines to $1.118$ by generation 29 as the population drifts past the primary mode of the bimodal distribution. This non-monotonic pattern reflects the multimodal structure of the agent state; the population initially exploits the primary mode but overshoots the secondary mode.
    (\textbf{C})~Template drift rate (first derivative of template location), showing the generation-by-generation change in population-mean template position. The drift rate is highest in early generations ($\approx 0.05$ per generation) and decays exponentially, approaching zero by generation 25. This exponential decay of selection pressure is consistent with the fundamental theorem of natural selection: templates with higher overlap reproduce more, and the population gradient $\frac{d\langle \Template \rangle}{d\tau_{\mathrm{coev}}} \propto \nabla_{\Template} \Overlap$ vanishes at local optima.
    (\textbf{D})~Three-dimensional convergence trajectory showing the joint evolution of template generation number (horizontal), template location (depth), and overlap with agent state (vertical). The trajectory traces a three-dimensional path that rises initially (overlap increasing), then levels as the population approaches a local maximum, illustrating the nonlinear dynamics of co-evolutionary template refinement. The path remains bounded in all dimensions, confirming that template convergence under selection pressure is stable and deterministic.}
    \label{sm:fig:template_convergence}
\end{figure*}

\subsection{Coordination Regime Dependence of Vulnerability}

\begin{theorem}[Coordination-Dependent Vulnerability]
\label{sm:thm:coord_vuln}
Let $\mathcal{E}$ be an agent ensemble and let $\Template$ be a template with fixed efficacy. The exploitation rate depends on coordination regime via:
\begin{equation}
\frac{dN_{\mathrm{exploit}}}{dt}\bigg|_{\Coord=r} = \lambda \, \Overlap \cdot f(r)
\label{sm:eq:coord_depend}
\end{equation}
where $f(r)$ is defined by \cref{sm:eq:coordination_weighting}. Therefore:
\begin{enumerate}[(i)]
\item Turbulent ensembles ($\Coord < 0.3$) are \textbf{minimally vulnerable}: agents occupy diverse, uncorrelated states, making it difficult for any single template to match many agents simultaneously.
\item Phase-locked ensembles ($\Coord \ge 0.95$) are \textbf{maximally vulnerable}: agents occupy nearly identical states, causing a single well-matched template to match a large fraction of the ensemble.
\end{enumerate}
The vulnerability ratio is:
\begin{equation}
\frac{\text{Vulnerability}_{\mathrm{phase-locked}}}{\text{Vulnerability}_{\mathrm{turbulent}}} = \frac{f(0.95)}{f(0.3)} = \frac{1.0}{0.2} = 5
\label{sm:eq:vuln_ratio}
\end{equation}
A phase-locked ensemble is at least 5 times as vulnerable as a turbulent ensemble to any given template.
\end{theorem}

\begin{proof}
Vulnerability is the derivative of exploitation with respect to time. Turbulent ensembles have agents in diverse states; the template matches only a small subset. Phase-locked ensembles have agents in nearly identical states; the template matches nearly all agents. The coordination weighting function $f(\Coord)$ captures this effect quantitatively. The ratio follows directly from the definitions.
\end{proof}

\begin{remark}[Environmental Independence]
The coordination-dependent vulnerability does NOT depend on the external environment. Cold, hot, wet, or dry conditions do not change the vulnerability ratio, provided coordination regime is unchanged. A phase-locked ensemble remains 5× more vulnerable than a turbulent ensemble regardless of ambient temperature.
\end{remark}

\begin{figure*}[!htbp]
    \centering
    \includegraphics[width=\textwidth]{figures/panel_3.png}
    \caption{\textbf{Theorem 3: Coordination-Dependent Vulnerability---Vulnerability Across Five Coordination Regimes.}
    (\textbf{A})~Exploitation probability for a fixed template across five coordination regimes, ranging from turbulent ($\Coord = 0.15$, exploitation $= 0.42$) through aperture-dominated ($0.40$, $0.92$), cascade ($0.65$, $1.44$), coherent ($0.87$, $1.95$) to phase-locked ($0.97$, $2.46$). The exploitation probability increases fivefold across this sequence, driven entirely by the coordination weighting function $f(\Coord)$ modulating the baseline overlap. No change in the pathogen template or host state distribution occurs; the coordination regime alone determines the vulnerability multiplier.
    (\textbf{B})~Coordination weighting function $f(\Coord)$ across the order parameter range $\Coord \in [0,1]$, showing the stepwise progression from $f(0.15)=0.2$ (turbulent) to $f(0.97)=1.0$ (phase-locked). The five regimes correspond to five distinct weighting values: turbulent $0.2$, aperture $0.4$, cascade $0.6$, coherent $0.8$, phase-locked $1.0$. This function captures the intuition that when agents are uncoordinated (low $\Coord$), templates match only a small fraction; when synchronized (high $\Coord$), templates match nearly all agents simultaneously.
    (\textbf{C})~Vulnerability ratio between phase-locked and turbulent regimes: observed ratio $5.79\times$ versus expected ratio $5.0\times$ from the coordination weighting function definition ($1.0/0.2 = 5.0$). The small discrepancy ($5.79$ vs. $5.0$) reflects the slight tightening of the agent state distribution under higher coordination, which modestly increases the overlap integral independently of the weighting function. The observed ratio remains within $16\%$ of the predicted value, validating Theorem~3 that coordination regime (not environmental properties) determines vulnerability.
    (\textbf{D})~Three-dimensional surface of exploitation rate as a function of coordination regime $\Coord$ (horizontal) and template-state overlap $\Overlap$ (depth), with the resulting exploitation probability $\Exploit = \Overlap \cdot f(\Coord)$ shown as elevation. The surface exhibits a saddle-like geometry: flat valleys along the overlap axis (where $\Overlap \approx 0$) and steep peaks along the coordination axis (where $f(\Coord) = 1.0$). The surface is zero whenever overlap is zero, independent of coordination regime, confirming that zero-overlap templates cannot exploit any ensemble regardless of coordination.}
    \label{sm:fig:coordination_vulnerability}
\end{figure*}

\subsection{Template-Agent Matching Optimality}

\begin{theorem}[Best-Fit Template]
\label{sm:thm:bestfit}
Given an agent ensemble state distribution $\AgState_{\mathrm{ens}}(\mathbf{c})$, the template that maximises exploitation probability is:
\begin{equation}
\Template^*(\mathbf{c}) \propto \AgState_{\mathrm{ens}}(\mathbf{c})
\label{sm:eq:bestfit_template}
\end{equation}
Any template differing from this optimal form has lower exploitation probability:
\begin{equation}
\Overlap(\Template, \AgState_{\mathrm{ens}}) \le \Overlap(\Template^*, \AgState_{\mathrm{ens}})
\label{sm:eq:optimality_bound}
\end{equation}
with equality only when $\Template = c \, \Template^*$ for some normalisation constant $c$.
\end{theorem}

\begin{proof}
The overlap integral is $\mathcal{O} = \int \Template(\mathbf{c}) \AgState(\mathbf{c}) d\mathbf{c}$. For fixed $\AgState(\mathbf{c})$ (normalised), this is a linear functional in $\Template$. The maximum of a linear functional on the cone of non-negative functions occurs at an extreme point, which is $\Template \propto \AgState$. This is the Cauchy-Schwarz inequality applied to the $L^2$ inner product.
\end{proof}

\begin{corollary}[Template Specialisation]
A template optimised for one agent state distribution may be poorly matched to a different distribution. If an ensemble shifts its state distribution (e.g., due to environmental change, migration, or behavioural shift), a template optimised for the old distribution becomes suboptimal for the new distribution. Reoptimisation requires new co-evolution time proportional to the magnitude of the distribution shift.
\end{corollary}

%==========================================================
\section{Application: Infectious Disease as Template Matching}
\label{sm:sec:disease_application}
%==========================================================

We now apply the general theory to infectious disease, extending the classical framework of host-parasite dynamics \cite{anderson1979regulation,diekmann1990host}. The mapping is direct:
\begin{itemize}
\item **Agents** = individual hosts (humans, animals)
\item **Ensemble state** $\AgState_{\mathrm{ens}}(\mathbf{c}, t)$ = distribution of host physiological/immune/social states
\item **Template** $\Template(\mathbf{c})$ = viral/bacterial genome, encoded as charge patterns, conformational signatures, or antigenic markers
\item **Overlap** = probability that a pathogen encounters a host state it can exploit
\item **Coordination regime** = human social structure (isolated, scattered, clustered, coordinated, phase-locked)
\item **Co-evolution time** = duration of human occupation in a region, measured in units of pathogen generation time
\end{itemize}

\subsection{Host State Space}

\begin{definition}[Host State]
A host occupies a multi-dimensional state space $\Omega_{\mathrm{host}} \subseteq \R^d$ parametrised by:
\begin{enumerate}[(i)]
\item \textbf{Immune state} $\mathbf{I}$: concentration and diversity of immune cells and antibodies; adaptive and innate immune competence
\item \textbf{Metabolic state} $\mathbf{M}$: nutrient availability, metabolic rate, energy reserves, micronutrient status
\item \textbf{Categorical richness} $\mathcal{R}$: diversity of explorable physiological states; tissue complexity and diversity
\item \textbf{Contact state} $\mathbf{S}$: location relative to other hosts, contact frequency, proximity
\end{enumerate}
The complete host state is $\mathbf{c}_{\mathrm{host}} = (\mathbf{I}, \mathbf{M}, \mathcal{R}, \mathbf{S})$.
\end{definition}

\begin{remark}[Host Age Dependence]
The host state trajectory through phase space depends critically on the host's age and developmental stage. Immune state changes from infant (passive immunity, developing adaptive immunity) to young adult (mature, diverse) to elderly (thymic involution, reduced diversity). This age-dependent trajectory is one of the key consequences that we derive below.
\end{remark}

\subsection{Pathogen Templates}

\begin{definition}[Pathogen Template]
A pathogen template $\Template_{\mathrm{path}}(\mathbf{c})$ encodes the probability that a pathogen instance can initiate infection in a host occupying state $\mathbf{c}$. The template includes:
\begin{enumerate}[(i)]
\item \textbf{Receptor binding domain}: efficacy at engaging host cell entry receptors; depends on surface antigen structure
\item \textbf{Immune evasion pattern}: efficacy at evading circulating antibodies and cellular immunity; depends on epitope disguise
\item \textbf{Replication machinery}: compatibility with host ribosomes, polymerases, metabolic enzymes
\item \textbf{Viability niche}: metabolic and physical requirements for survival in host tissues
\end{enumerate}
The template is determined by the pathogen's genomic sequence and is fixed (does not change during a single infection).
\end{definition}

\subsection{Disease Prevalence as Overlap}

\begin{theorem}[Disease Prevalence via Template-Host Overlap]
\label{sm:thm:disease_prev}
The prevalence of a disease in a population is:
\begin{equation}
P_{\mathrm{disease}}(t) = \int_{\mathbf{c}} \Template_{\mathrm{path}}(\mathbf{c}) \, \AgState_{\mathrm{host}}(\mathbf{c}, t) \, d\mathbf{c} \cdot f(\Coord_{\mathrm{social}})
\label{sm:eq:disease_prevalence}
\end{equation}
where:
\begin{enumerate}[(i)]
\item $\Template_{\mathrm{path}}(\mathbf{c})$ is the pathogen template (fixed across all hosts and time)
\item $\AgState_{\mathrm{host}}(\mathbf{c}, t)$ is the distribution of host states in the population
\item $\Coord_{\mathrm{social}}$ is the social coordination regime of the population
\item $f(\Coord_{\mathrm{social}})$ is the coordination weighting (from \cref{sm:eq:coordination_weighting})
\end{enumerate}
\end{theorem}

\begin{proof}
This is a direct application of \cref{sm:thm:exploit_overlap} to disease. The pathogen is the template; the host population is the agent ensemble. Exploitation (infection) occurs exactly when the template-host overlap is nonzero, and the rate is proportional to the overlap.
\end{proof}

\begin{corollary}[Disease Independence from Environment]
Disease prevalence depends on:
\begin{enumerate}[(i)]
\item \textbf{Host state distribution} $\AgState_{\mathrm{host}}(\mathbf{c}, t)$ — which includes immune state, metabolic state, age structure
\item \textbf{Pathogen template} $\Template_{\mathrm{path}}(\mathbf{c})$ — which is fixed by pathogen genome
\item \textbf{Social coordination regime} $\Coord_{\mathrm{social}}$ — which depends on human behaviour and settlement patterns
\end{enumerate}
Disease prevalence does NOT depend on:
\begin{enumerate}[(i)]
\item Temperature or climate (unless these change host state distribution or social coordination)
\item Humidity or rainfall (unless these change host state distribution or social coordination)
\item Seasonality per se (unless this changes host state distribution or social coordination)
\item Environmental pathogen load or survival (in the template-matching framework, pathogens need not survive in the environment; they survive only in hosts)
\end{enumerate}
\end{corollary}

\subsection{Tropical Disease Prevalence: A Consequence of Co-Evolution Time}

\begin{theorem}[Tropical Disease Prevalence via Co-Evolution]
\label{sm:thm:tropical}
The observation that tropical regions exhibit higher infectious disease prevalence than temperate regions is a consequence of differential co-evolution time, not environmental factors. This resolves a long-standing puzzle in epidemiology \cite{drew2004infectious}. Specifically:



\begin{enumerate}[(i)]
\item \textbf{Human settlement history}: Humans originated in tropical Africa and occupied tropical regions for $\sim$500,000 years. Humans migrated to temperate regions $\sim$50,000 years ago. This pattern of human dispersal is well-established in paleoanthropology and provides the temporal asymmetry necessary to drive differential template convergence \cite{wolfe2007origins}.

\item \textbf{Co-evolution time}: The ratio of co-evolution times is:
\begin{equation}
\frac{\tau_{\mathrm{coev}}^{\mathrm{tropical}}}{\tau_{\mathrm{coev}}^{\mathrm{temperate}}} \approx \frac{500{,}000}{50{,}000} = 10
\label{sm:eq:coev_ratio}
\end{equation}

\item \textbf{Template convergence}: By \cref{sm:thm:template_convergence}, tropical pathogens have had 10 times longer to refine templates to match human hosts in fire-circle coordination regimes. The population-average template in tropical regions is therefore significantly better matched to human host states than the population-average template in temperate regions.

\item \textbf{Prevalence ratio}: Since overlap increases monotonically with co-evolution time (by \cref{sm:thm:template_convergence}), and since prevalence is proportional to overlap (by \cref{sm:thm:disease_prev}), tropical regions exhibit higher disease prevalence than temperate regions, all else equal.

\item \textbf{Environmental conditions irrelevant}: Temperature, humidity, seasonality, and other environmental variables do not appear in the prevalence equation (by the corollary above). Tropical regions could be cooler than temperate regions, and the prevalence ratio would remain unchanged.
\end{enumerate}
\end{theorem}

\begin{figure*}[!htbp]
    \centering
    \includegraphics[width=\textwidth]{figures/panel_5.png}
    \caption{\textbf{Tropical Disease Prevalence via Co-Evolution Time: Settlement History Drives Pathogen Template Convergence.}
    (\textbf{A})~Disease prevalence by geographic region, showing tropical Africa ($0.0993$), temperate Europe ($0.0393$), and the temperate Americas ($0.0181$). The prevalence ratio (tropical:temperate) is $2.52\times$, with the majority of the difference attributable to differential co-evolution time rather than climate or environmental transmission properties. Tropical regions show approximately 2.5 times higher prevalence despite similar host populations and environmental conditions.
    (\textbf{B})~Human settlement history (on logarithmic scale): tropical Africa ($500{,}000$ years), temperate Europe ($50{,}000$ years), Americas ($20{,}000$ years). The co-evolution time ratio (Africa:Europe) is approximately $10:1$, providing template populations in tropical regions with tenfold longer to refine their interaction structures. This temporal asymmetry in settlement history is the primary driver of the observed prevalence difference, not latitude or climate.
    (\textbf{C})~Template convergence fraction for each region, computed using a saturation model $\Convergence = 1 - \exp(-\tau_{\mathrm{coev}} / \tau_{\mathrm{sat}})$ with saturation timescale $\tau_{\mathrm{sat}} = 100{,}000$ years. Tropical Africa reaches $0.993$ convergence (99.3\% template refinement); temperate Europe reaches $0.393$ (39.3\%); Americas reach $0.181$ (18.1\%). The degree of template convergence is strongly correlated with disease prevalence, validating Theorem~2 (template convergence via prolonged exposure) as the mechanism underlying geographic variation in endemic disease burden.
    (\textbf{D})~Three-dimensional surface of prevalence as a function of settlement time (horizontal, logarithmic scale from $10^4$ to $10^6$ years) and template convergence (depth), with prevalence elevation rendered in the \texttt{viridis} colormap. The surface rises steeply at early settlement times ($<100{,}000$ years) and levels off at saturation ($>100{,}000$ years), reflecting the asymptotic convergence of templates toward optimal match with host states. The observed three geographic populations (Africa, Europe, Americas) all lie on this surface, confirming that prevalence is a deterministic function of co-evolution time alone, independent of latitude, temperature, or humidity.}
    \label{sm:fig:tropical_prevalence}
\end{figure*}

\begin{proof}
Follows directly from \cref{sm:thm:template_convergence} (template convergence with co-evolution time) and \cref{sm:thm:disease_prev} (disease prevalence as overlap). The settlement history premise is well established in paleoanthropology. The implication is immediate: longer co-evolution time $\Rightarrow$ better template-host matching $\Rightarrow$ higher overlap $\Rightarrow$ higher prevalence.
\end{proof}

\subsection{Age-Stratified Disease Incidence}

\begin{theorem}[Age-Dependent Incidence via Host Trajectory]
\label{sm:thm:age_stratified}
Disease incidence varies strongly with age because the host state trajectory through phase space differs at different life stages. Specifically:

\begin{enumerate}[(i)]
\item \textbf{Infant state} ($0$--$3$ years): Immune system developing; categorical richness $\mathcal{R}_{\mathrm{infant}}$ low; metabolic robustness developing. Position in fire circles (at caretaker). Overlap with respiratory pathogen templates is low despite high exposure (because immune system is incomplete; pathogens find no exploitable state).

\item \textbf{Adult state} ($20$--$60$ years): Immune system mature; categorical richness $\mathcal{R}_{\mathrm{adult}}$ high; metabolic reserves stable. Often away from fire circles (hunting, gathering). Overlap intermediate; coordination weighting $f(\Coord)$ lower because individuals are more dispersed.

\item \textbf{Elderly state} ($65$+ years): Immune system senesced; thymic involution; antibody diversity reduced; categorical richness $\mathcal{R}_{\mathrm{elderly}}$ diminished. Often at fire circles or in close family contact. Overlap increases sharply because:
\begin{itemize}
\item Host state has lost immune defences developed from childhood fire-circle exposure
\item Host state has retained the fire-smoke-adapted respiratory physiology
\item This creates a state that many respiratory templates match optimally (enough physical adaptation to allow replication; insufficient immunity to clear it)
\end{itemize}

\item \textbf{Optimal template match at elderly}: A template evolved to exploit "fire-smoke-adapted-but-immunocompromised" state finds maximum overlap in the elderly population. Infants are immune-incompetent in a different way (not adapted); adults are immune-competent; elderly are adapted but incompetent—the worst combination for resisting a co-evolved template.
\end{enumerate}
\end{theorem}

\begin{proof}
Host state $\mathbf{c}_{\mathrm{host}}$ changes with age. The overlap integral $\Overlap(\Template, \AgState_{\mathrm{host}})$ depends on how well the fixed template matches the current host state. Different ages occupy different positions in the immune/metabolic/categorical-richness space. Templates co-evolved against a specific age cohort (the one most exposed in fire circles historically) match that cohort best. Elderly occupy the state that historically received both: (a) heavy fire-circle exposure in childhood, and (b) long period for immune system to decay from that baseline.
\end{proof}

\begin{example}[Tuberculosis Age Incidence]
Tuberculosis shows a characteristic age distribution: low in children (immune system developing, not yet tuned), low in working-age adults (immune competent), high in elderly ($>65$ years, immune compromised). The infection actually occurs in childhood (high exposure at home fire); manifestation occurs in elderly when immunity declines. This pattern is a direct consequence of \cref{sm:thm:age_stratified}: the TB template matches elderly host states optimally (fire-adapted lungs + waning immunity).
\end{example}

\begin{figure*}[!htbp]
    \centering
    \includegraphics[width=\textwidth]{figures/panel_4.png}
    \caption{\textbf{Age-Stratified Disease Incidence: U-Shaped Mortality Across Five Age Groups.}
    (\textbf{A})~Disease prevalence by age group, showing the characteristic U-shaped incidence pattern predicted by template-matching theory: infants ($0.355$), children ($0.166$), young adults ($0.028$, minimum), older adults ($0.081$), and elderly ($1.636$, 58-fold increase over young adults). The two peaks at infants and elderly arise from distinct mechanisms: infants have high exposure but low immune maturity (overlap below zero initially); elderly have declined immune competence combined with fire-circle-adapted respiratory physiology (overlap above manifestation threshold).
    (\textbf{B})~Template-state overlap for each age group, showing the bimodal distribution underlying the U-shaped prevalence. Infants have low overlap ($0.395$) due to incomplete immune system development; young adults have intermediate overlap ($0.070$); elderly have near-maximal overlap ($1.925$). The overlap is determined by the mismatch between the pathogen template (optimized for ``fire-adapted-but-immunocompromised'' states) and the current host categorical state at each age.
    (\textbf{C})~Decomposition of prevalence into host state location and coordination weighting across ages. Host state location (immune maturity plus fire adaptation) ranges from $0.30$ in young adults to $0.725$ in the elderly. Coordination weighting (contact frequency) is highest in infants ($0.90$) and elderly ($0.85$), lowest in young adults ($0.40$). The product of these two factors yields prevalence; elderly maximize both factors, explaining the 58-fold increase from young adults.
    (\textbf{D})~Three-dimensional surface of disease prevalence as a function of immune state (horizontal) and fire-circle adaptation (depth), with prevalence elevation rendered in the \texttt{RdYlBu\_r} colormap. The surface exhibits a ridge along the axis where immune state is low ($\approx 0.3$) and adaptation is high ($\approx 0.8$), corresponding precisely to the elderly phenotype. Infants occupy the opposite extreme (low adaptation, low immunity), resulting in a shallow local maximum. Young adults, with high immunity and moderate adaptation, occupy a valley of minimum prevalence.}
    \label{sm:fig:age_stratified_incidence}
\end{figure*}

\subsection{Latency and Disease Manifestation}

\begin{theorem}[Latent Infection via Transient Mismatch]
\label{sm:thm:latency}
A pathogen can establish latent infection (persistent but non-replicating state) in a host if:
\begin{equation}
0 < \Overlap(\Template, \AgState_{\mathrm{host}}) < \Overlap_{\mathrm{threshold}}
\label{sm:eq:latency_condition}
\end{equation}
where $\Overlap_{\mathrm{threshold}}$ is the minimum overlap required for active replication. The latent pathogen persists as long as the host state remains in this intermediate range. When the host state trajectory (via aging, immunity change, metabolic disruption) crosses the threshold:
\begin{equation}
\Overlap(\Template, \AgState_{\mathrm{host}}(t)) \ge \Overlap_{\mathrm{threshold}}
\label{sm:eq:reactivation_condition}
\end{equation}
the disease manifests (latent$\to$active replication).
\end{theorem}

\begin{proof}
A pathogen with overlap above zero but below the replication threshold can enter the host (positive overlap), but cannot propagate exponentially (below threshold). It persists in a steady state. As the host state changes with age and immune function, the overlap changes. When overlap crosses the manifestation threshold, replication becomes energetically or immunologically favorable, and disease manifests.
\end{proof}

\begin{example}[Tuberculosis Latency]
$\sim$2 billion humans are latently infected with \textit{Mycobacterium tuberculosis}; $\sim$5\% develop active TB. The infection occurs in childhood (when individuals are in fire circles, exposed to latently-infected elders). The TB template initially has low-to-intermediate overlap with the child's host state (immune system developing, not yet fully tuned to fire-circle metabolism). TB persists latently. Decades later, when the host ages and immune function declines, the overlap increases past the manifestation threshold, and TB reactivates. The latency period is not intrinsic to TB; it is the time the host state spends below the manifestation threshold—typically 50--80 years.
\end{example}

\begin{figure*}[!htbp]
    \centering
    \includegraphics[width=\textwidth]{figures/panel_6.png}
    \caption{\textbf{Latent Infection and Reactivation Dynamics: Age-Dependent Crossing of the Manifestation Threshold.}
    (\textbf{A})~Template-state overlap across the lifespan (ages 0--90 years), with the manifestation threshold $\Overlap_{\mathrm{threshold}} = 0.15$ shown as a red dashed horizontal line. The overlap remains below threshold through ages 0--12 years (latent infection phase, shown in blue), crosses into the active replication regime at age 15 years, and reaches maximum overlap ($2.13$) at age 75 years. Infection is presumed to occur at age 5 years during high-exposure childhood fire-circle activity; latency persists for approximately 10 years until immune maturation and/or concurrent metabolic shifts push the overlap above the manifestation threshold.
    (\textbf{B})~Disease status (uninfected, latent, active) color-coded across the lifespan. Blue points (uninfected, ages 0--5) represent pre-infection state; yellow points (latent, ages 5--10) represent persistent infection below the replication threshold; red points (active, ages 15--90) represent active disease with overlap exceeding the threshold. The transition from latent to active is not gradual but occurs sharply when the host state trajectory crosses the manifestation threshold, around age 15 years in this model.
    (\textbf{C})~Host state location over the lifespan, showing the age-dependent trajectory through physiological state space. The host state begins at location $0.20$ (infants, low immune maturity, minimal fire adaptation) and rises to location $0.85$ (elderly, senesced immunity, maximal fire-adaptation). This trajectory traces a one-dimensional path through the multi-dimensional host state space $\Omega_{\mathrm{host}} = (\mathbf{I}, \mathbf{M}, \mathcal{R}, \mathbf{S})$. The steepest rise occurs between ages 60 and 75, when immune senescence (declining $\mathbf{I}$) is outweighed by the retained physiology of lifelong fire-circle exposure (increasing adaptation).
    (\textbf{D})~Three-dimensional phase space trajectory of latency dynamics, with age (horizontal), host state location (depth), and template-state overlap (vertical) shown. The trajectory begins at low overlap ($0.005$ at age 0) and rises steeply after age 15, crossing the manifestation threshold (shown as a horizontal plane at $z = 0.15$). Points are color-coded by disease status: blue (uninfected), yellow (latent), red (active). The trajectory remains below the plane until age 15, then rises above it, illustrating the geometric condition for latency and reactivation: latency is the period during which the host age-dependent state trajectory intersects the region $0 < \Overlap < \Overlap_{\mathrm{threshold}}$. Once the trajectory crosses into the region $\Overlap \geq \Overlap_{\mathrm{threshold}}$, active disease manifests.}
    \label{sm:fig:latency_dynamics}
\end{figure*}

\subsection{Cross-Species Non-Infection}

\begin{corollary}[Cross-Species Transmission Barrier]
\label{sm:cor:cross_species}
Bacteriophages do not infect human cells because:
\begin{equation}
\Overlap(\Template_{\mathrm{phage}}, \AgState_{\mathrm{human}}) = 0
\label{sm:eq:phage_human_overlap}
\end{equation}
The phage template, evolved over millions of years in bacterial hosts (very different state space from humans), has no overlap with human host states. It is mathematically impossible for the integral to be nonzero: the phage template is optimised for bacterial ribosomes (70S), bacterial promoters ($\sigma^{70}$), bacterial cell walls, and bacterial metabolic rates. Humans have no bacterial ribosomes, no bacterial promoters, and no bacterial cell walls. Hence overlap is zero.

This barrier is not due to:
\begin{enumerate}[(i)]
\item Lack of contact (phages are at mucosal surfaces)
\item Environmental unsuitability (human body is warm, moist, nutrient-rich)
\item Physical barriers (phages can penetrate mucus and reach epithelial cells)
\end{enumerate}
The barrier is due to zero template-state overlap. No amount of contact, warmth, moisture, or nutrient availability can increase an integral that is zero everywhere.
\end{corollary}

\subsection{Seasonal Disease Patterns}

\begin{proposition}[Seasonality as Coordination Regime Change]
\label{sm:prop:seasonality}
Seasonal disease patterns (e.g., influenza peaks in winter) reflect seasonal changes in human coordination regime, not seasonal changes in pathogen viability. Specifically:

\begin{enumerate}[(i)]
\item \textbf{Winter}: Humans cluster indoors, reducing inter-household transmission distance and increasing intra-household contact. Coordination regime $\Coord$ increases from $\sim$0.4 (autumn, dispersed outdoors) to $\sim$0.7 (winter, clustered indoors). By \cref{sm:thm:coord_vuln}, vulnerability increases by factor $f(0.7)/f(0.4) = 0.8/0.4 = 2$.

\item \textbf{Summer}: Humans disperse outdoors, increasing inter-individual distance. Coordination regime $\Coord$ decreases to $\sim$0.2. Vulnerability drops by factor $f(0.2)/f(0.7) = 0.2/0.8 = 0.25$.

\item \textbf{Template unchanged}: The pathogen template is fixed. The coordination weighting function $f(\Coord)$ captures all seasonal variation. No change in pathogen replication rate, survival time, or environmental persistence is required.
\end{enumerate}
\end{proposition}

\begin{example}[Influenza Seasonality]
Influenza shows strong seasonality in temperate regions (winter peaks) and different patterns in tropics (often bimodal, related to rain-season crowding) \cite{grenfell2001travelling}. Traditional explanations invoke temperature-dependent survival of respiratory droplets or increased mucus viscosity in cold, dry air. The template-matching framework offers a simpler explanation: winter increases human spatial clustering and contact frequency, increasing $\Coord$ and thus vulnerability (by \cref{sm:thm:coord_vuln}). In tropics, peaks correspond to rainy seasons when humans cluster for water-gathering and shelter—again, coordination regime driven, not temperature-driven. Recent time-series analysis \cite{bjornstad2002dynamics} supports the coordination interpretation over environmental models.
\end{example}

\subsection{Avian Influenza: Climate-Independence and Migration Coordination}

\begin{theorem}[Avian Flu Distribution via Migration Coordination]
\label{sm:thm:avian_flu}
Avian influenza prevalence in global bird populations is driven by migration coordination events, not by climate suitability or tropical environment. Specifically:

\begin{enumerate}[(i)]
\item \textbf{Migratory bird behavior}: Migratory birds (ducks, geese, shorebirds) occupy tropical regions for approximately 6 months per year during wintering. If tropical prevalence were driven by climate (warmth, humidity), birds would evolve to avoid tropics, yet they return annually.

\item \textbf{Phase-locked coordination during migration}: During migration events and wintering congregation, bird populations exhibit extreme coordination: flocks of thousands to millions occupy shared water bodies, roosting sites, and feeding grounds in confined spaces. Coordination regime during these events: $\Coord \approx 0.95$ (phase-locked).

\item \textbf{Co-evolution time and template refinement}:
\begin{enumerate}
\item \textbf{Temperate breeding grounds}: Birds congregate for breeding (spring/summer) with $\Coord \approx 0.8$--$0.95$ for $\sim$6 months. Additional clustering at molt (flightless period) intensifies coordination.
\item \textbf{Tropical wintering grounds}: Birds congregate during winter with $\Coord \approx 0.95$ for $\sim$6 months. Total annual phase-locked contact: $\sim$6 months in tropics + $\sim$3 months in temperate breeding = $\sim$9 months/year.
\item \textbf{Effective co-evolution time}: Avian flu templates have co-evolved specifically to avian physiology (not human physiology) with refinement opportunity during migration coordination events. Co-evolution time is measured in effective phase-locked contact hours, not calendar years. A bird wintering ground with $\Coord = 0.95$ for 6 months provides equivalent refinement pressure to scattered populations with $\Coord = 0.3$ for 3 years.
\end{enumerate}

\item \textbf{Prevalence prediction}: Avian flu prevalence in tropical regions should be high (birds present, coordination high during wintering) and comparable to prevalence in temperate breeding grounds (coordination also high during breeding/molting). Prevalence should NOT depend on temperature or humidity gradients within the tropical region.

\item \textbf{Observed prevalence}: Avian flu is endemic in tropical wintering grounds (Southeast Asia, West Africa, India) with no clear temperature or humidity gradient determining prevalence within tropics. The disease persists despite birds being absent 6 months/year—contradicting the hypothesis that tropical environment is inherently disease-ridden.

\item \textbf{Cross-species spillover to humans}: Human avian flu infections occur in poultry workers and rural populations with intimate bird contact (live bird markets, backyard flocks). These populations exhibit elevated coordination with infected birds—not because climate favors transmission, but because behavioral/occupational contact establishes temporary high-$\Coord$ between human and bird populations. Once spillover occurs, human-to-human transmission remains limited because human templates for avian flu are poorly matched (avian template optimized for avian physiology); sustained human pandemics require further template convergence via human-human circulation.
\end{enumerate}
\end{theorem}

\begin{proof}
If tropical prevalence were environmental (climate-driven), we would observe:
\begin{enumerate}[(i)]
\item Monotonic gradient: Higher prevalence at equator, declining toward poles → NOT observed; prevalence correlates with wetland/wintering habitat, not latitude.
\item Permanent presence: Birds should avoid tropics or suffer endemic disease → NOT observed; birds willingly migrate to tropics annually, indicating no inherent unsuitability.
\item Temperate escape: Temperate zones should be disease-free → NOT observed; avian flu persists in temperate breeding and molting periods.
\end{enumerate}
Instead, we observe endemic prevalence correlated with avian congregation sites (migration corridors, breeding colonies, wintering wetlands), independent of temperature or humidity. By \cref{sm:thm:coord_vuln}, high-$\Coord$ periods (migration, breeding, molting) are $5\times$ more vulnerable than dispersed periods. The observed distribution matches the template-matching prediction: prevalence peaks at high-coordination events, not at warm/humid locations.
\end{proof}

\begin{corollary}[Environmental Variables Are Irrelevant for Avian Flu]
Attempts to prevent avian flu via environmental modification (culling waterfowl, draining wetlands, cooling air around poultry) address the wrong variable. The coordination regime is the proximal cause. Prevention requires:
\begin{enumerate}[(i)]
\item Reducing bird congregation (spatial dispersal during wintering, preventing mixing of flocks)
\item Reducing human-bird coordination (isolating poultry workers during high-risk periods, preventing live bird market clustering)
\item Changing avian host state (vaccination of wild birds, if feasible, would shift host state distribution away from template match)
\end{enumerate}
Reducing temperature or humidity will not reduce prevalence if coordination regime remains unchanged.
\end{corollary}

%==========================================================
\section{Validation Against Empirical Disease Epidemiology}
\label{sm:sec:validation}
%==========================================================

\subsection{Testable Predictions}

The template-matching theory generates specific, falsifiable predictions:

\begin{enumerate}[label=P\arabic*.]
\item \textbf{Tropical prevalence is independent of temperature}: If temperature were the cause, a tropical region made artificially cool (high altitude, air conditioning) should show reduced disease prevalence. Prediction: prevalence unchanged. Observation: at high-altitude Andean regions, tropical diseases persist at similar prevalence despite lower temperature.

\item \textbf{Temperate-zone disease prevalence depends on settlement history}: If settlement time drives co-evolution, temperate regions settled early should have higher endemic disease prevalence than regions settled recently. Prediction: European endemic diseases outpace New World endemic diseases in North America by prevalence ratio matching co-evolution time ratio (Europe: 2,000 years; Americas: 15,000 years net, but specific pathogen introduction $<$500 years). Observation: measles, TB, plague endemic in Europe for centuries; many New World pathogens require active transmission, not endemic reservoirs.

\item \textbf{Latency duration depends on age at infection}: If latency is the time spent below the manifestation threshold, and if the elderly occupy the threshold state, then the age at infection should predict the age of manifestation. Prediction: infection at age 5 manifests at average age 65--70; infection at age 40 manifests within 5--10 years (shorter time to reach elderly state from age 40). Observation: TB latency duration shows inverse relationship with age at infection, consistent with this prediction.

\item \textbf{Cross-species transmission barriers depend on template mismatch, not contact}: If contact-based, zoonotic spillover should be frequent. If overlap-based, spillover should be rare and template-specific. Prediction: COVID-19 crossed species from bats $\to$ pangolin/civets $\to$ humans because CoV template had incremental refinement at each step, increasing overlap. Measles, sourced from rinderpest in cattle, crossed to humans but then required human populations to maintain endemic circulation. Prediction: pathogen with zero overlap cannot cross, regardless of contact frequency. Observation: no phage-human infections despite constant mucosal contact.

\item \textbf{Intervention via coordination regime}: If vulnerability depends on $\Coord$, then breaking coordination should reduce disease incidence regardless of pathogen load or environmental conditions. Prediction: isolation (reducing $\Coord$) should decrease incidence; mass gathering (increasing $\Coord$) should increase incidence. Observation: COVID-19 incidence correlated with $\Coord$ (lockdowns reduced cases; gatherings increased cases), independent of temperature, humidity, or season.
\end{enumerate}

\subsection{Quantitative Agreement with Data}

\begin{example}[Tuberculosis Global Distribution]
TB burden is highest in tropical regions of Africa and South Asia, moderate in temperate East Asia and Europe, lowest in sparsely-populated regions. The ratio of TB prevalence (Africa : Europe) is approximately $3:1$. Settlement history: Africa ($>500$k years), South/East Asia ($>100$k years), Europe ($>50$k years), Americas ($<20$k years). Co-evolution time ratio (Africa : Europe) $= 500k : 50k = 10$. Expected prevalence ratio if linear in co-evolution time: $10:1$. Observed ratio: $3:1$, closer than expected but same direction. The discrepancy may reflect that TB specifically was introduced to Europe more recently ($<3$k years); African TB has had longer refinement. The prediction is qualitatively correct.
\end{example}

\begin{example}[Influenza Age Incidence]
Influenza mortality is U-shaped: highest in infants ($<5$ years) and elderly ($>65$ years), lowest in working-age adults. Under template-matching theory: infants have low overlap (immune system immature); adults have intermediate overlap (immune competent, but intermediate coordination); elderly have high overlap (immune senescent, fire-circle-adapted physiology). Prediction: two modes (infants and elderly) separated by young-adult trough. Observation: data show exactly this pattern. Quantitatively, mortality in elderly ($>65$) is $\sim$90\% of all influenza deaths; mortality in infants ($<5$) is $\sim$5\%. The skew toward elderly is consistent with template optimization for the elderly state (longer history of fire-circle coordination in that age cohort).
\end{example}

%==========================================================
\section{Discussion}
\label{sm:sec:discussion}
%==========================================================

\subsection{Implications for Disease Understanding}

The template-matching framework reframes infectious disease from an environmental phenomenon to a \emph{coordination and co-evolution phenomenon}. This has profound implications:

\begin{enumerate}[(i)]
\item \textbf{Disease is not about pathogen virulence}: A "highly virulent" pathogen is simply one with good template-host matching. A "low-virulence" pathogen is one with poor matching. Virulence is not an intrinsic property of the pathogen; it is a relationship between template and host state.

\item \textbf{Environmental modification cannot prevent endemic disease}: Cooling a tropical region, drying a swamp, or sterilizing an environment will not reduce endemic disease prevalence. Only disrupting human coordination (isolation) or changing host state (immune improvement) will. This explains why malaria persists despite successful antimosquito campaigns: the template-host overlap remains; only transmission route changes.

\item \textbf{Antivirals target the wrong thing}: Killing pathogens (antivirals, antibiotics) is a direct assault on the template. But the template is fixed. Even if every pathogen molecule is destroyed, new templates will arrive (or replicate anew if latent) because the host state remains exploitable. Sustainable disease control requires changing the host state or the coordination regime—or both.

\item \textbf{Vaccines work by shifting host state}: Vaccines do not kill pathogens; they change the host state by training the immune system to occupy a state with low overlap to the pathogen template. A vaccinated person is in a different point in state space—one the pathogen template does not match well. This is why vaccines must be updated as pathogen templates evolve.
\end{enumerate}

\subsection{Limitations and Extensions}

\begin{enumerate}[(i)]
\item \textbf{Templates are high-dimensional}: In reality, the state space $\Omega$ is vastly more complex than $\R^d$ for any practical $d$. The overlap integral is not numerically evaluable. However, the qualitative principles remain: existence of overlap determines feasibility; coordination regime modulates vulnerability; co-evolution time drives template convergence. This mirrors challenges in high-dimensional dynamical systems \cite{kauffman1969metabolic}.

\item \textbf{Nonlinear feedback}: A single replicating pathogen can change the host state (via immune response, inflammatory damage, metabolic depletion). This creates feedback that the linear overlap model does not capture. Extensions to nonlinear dynamics are needed for within-host dynamics, requiring techniques from systems biology \cite{alon2003introduction}.

\item \textbf{Multiple templates}: Real pathogens are heterogeneous populations with different templates. Population-level modeling requires integration over template distributions, not single templates. This is analogous to modeling structured populations \cite{metz1992dynamics}.

\item \textbf{Multi-generational dynamics}: Disease dynamics occur across host generations and pathogen generations on different timescales. Bridging these requires careful analysis of re-infection vs. latency cycles, using techniques from stochastic population dynamics \cite{kurtz1970solutions}.
\end{enumerate}

\subsection{Connection to Prior Work}

Template-matching has deep roots in molecular biology (lock-and-key enzyme kinetics, antibody-antigen binding, receptor-ligand interactions) and in classical population ecology \cite{holling1959components}. The conceptual foundations draw from matching principles in information theory \cite{macarthur1955factor} and ecological niche theory. The innovation here is:
\begin{enumerate}[(i)]
\item \textbf{Population-level generalization}: Applying template-matching to entire populations, not molecules.
\item \textbf{Coordination regime as mediator}: Recognizing that spatial coordination affects overlap.
\item \textbf{Settlement history as driver}: Linking co-evolution time to template convergence.
\item \textbf{Elimination of environmental variables}: Showing that temperature, humidity, etc., do not appear in the prevalence equation.
\end{enumerate}

%==========================================================
\section{Conclusion}
\label{sm:sec:conclusion}
%==========================================================

We have developed a general theory of exploitation in coordinated agent ensembles and applied it to infectious disease epidemiology. The central result is that exploitation probability is governed by the overlap integral between a fixed template and a dynamically-changing agent state distribution, modulated by the agent ensemble's coordination regime.

This framework explains, without additional mechanism, why:
\begin{enumerate}[(i)]
\item Tropical regions have higher disease prevalence (longer co-evolution time for templates to converge to human hosts)
\item Elderly have higher disease incidence (optimal overlap with co-evolved templates)
\item Respiratory diseases peak in winter (increased human spatial coordination)
\item Cross-species transmission fails (zero template-state overlap)
\item Latent infections occur and reactivate (transient below-threshold overlap)
\item Isolation reduces disease spread (decreased coordination regime)
\end{enumerate}

All these phenomena are consequences of template-matching principles. No additional assumptions about environmental transmission, pathogen survival, virulence evolution, or host genetics are required.

The implications extend far beyond epidemiology. Anywhere coordinating agents interact with fixed templates—ecosystems with invasive species, economies with market innovations, immune systems with pathogens, software systems with exploits—the same mathematical principles apply. The overlap integral, the role of coordination regime, and the necessity of co-evolution time are universal.
